\documentclass[11pt]{article}

\usepackage[a4paper,margin=2.5cm]{geometry}
\usepackage{amsmath,amssymb,amsthm}
\usepackage[colorlinks=true,linkcolor=blue,citecolor=blue,urlcolor=blue]{hyperref}
\usepackage{booktabs}
\usepackage{enumitem}
\usepackage{natbib}

\newtheorem{theorem}{Theorem}[section]
\newtheorem{lemma}{Lemma}[section]
\newtheorem{proposition}{Proposition}[section]
\newtheorem{definition}{Definition}[section]
\newtheorem{remark}{Remark}[section]
\newtheorem{hypothesis}{Hypothesis}[section]

\newcommand{\twoI}{2I}
\newcommand{\Vsixhundred}{V_{600}}
\newcommand{\Dicfive}{\mathrm{Dic}_5}
\newcommand{\CQ}{C_{\mathbb{Q}}}
\newcommand{\IC}{I_{\CQ}}
\newcommand{\PK}[1]{P_{#1}}
\newcommand{\Hop}{\hat H}
\newcommand{\Cop}{\hat C}
\newcommand{\Ttau}{T_{\tau}}
\newcommand{\Hzero}{H_0}
\newcommand{\sigeight}{\sigma_8}

\title{Two operator-trace ratios from $\Vsixhundred$ K-class structure:\\
       $13/12$ and $11/12$, and their match to the
       $\Hzero$ and $S_8$ benchmarks}
\author{%
  Lee Smart \\[3pt]
  \textit{Institute of Vibrational Field Dynamics} \\[2pt]
  \texttt{contact@vibrationalfielddynamics.org}}
\date{May 2026}

\begin{document}

\maketitle

\noindent\textbf{Status:} Preprint (not peer-reviewed).
This paper has a deliberate three-layer structure: Layer~1 is exact
mathematics, Layer~2 is a numerical observation from published data,
Layer~3 is an explicit hypothesis. We do NOT claim Layer~2 or Layer~3
are theorems.

\begin{abstract}
We define two operators on the $12$-dimensional rational
$\Ttau$-cycle space $\CQ \cong \mathbb{Q}^{12}$ of the $600$-cell:
$\Hop := \IC + \PK{72}$ and $\Cop := \IC - \PK{0}$, where $\PK{K}$ is
the diagonal rank-$m_K$ projector onto the K-class cycles in the
T$_\tau$-cycle decomposition with K-multiset
$\{72{:}1,\,0{:}1,\,52{:}5,\,20{:}5\}$
(Paper~1, \citealp{V600Paper1}). We prove (Theorem~\ref{thm:trace})
the exact rational identities $\mathrm{tr}(\Hop) = 13$ and
$\mathrm{tr}(\Cop) = 11$, and we prove
(Theorem~\ref{thm:admissibility}) that within the K-saturated
diagonal projector algebra of $\CQ$ the only rank-one projector corrections to
$\IC$ are $\pm\PK{72}$ and $\pm\PK{0}$. Both theorems are
phase-independent across the $\mathbb{Z}_2^5 = 32$ canonical lifts of
$\tau_\sigma$ (Paper~3, \citealp{V600Paper3}).

As a numerical observation, we then report:
$\frac{13}{12} \cdot \Hzero^{\mathrm{Planck}} \approx 72.97\,
\mathrm{km\,s^{-1}\,Mpc^{-1}}$ versus
SH0ES~$73.04 \pm 1.04$~\citep{Riess2022SH0ES}
($\approx 0.06\sigma$ residual), and
$\frac{11}{12} \cdot S_8^{\mathrm{Planck}} \approx 0.763$ versus
KiDS-1000~$0.766^{+0.020}_{-0.014}$~\citep{Heymans2021KiDS}
($\approx 0.24\sigma$ directional residual using the lower KiDS
asymmetric error; we use the lensing parameter
$S_8 = \sigeight \sqrt{\Omega_m / 0.3}$ since published weak-lensing
intervals are reported as $S_8$).
We do NOT derive cosmology from operator traces; the cycle-mode
coupling rule that would upgrade the numerical match to a derivation
is stated in \S\ref{sec:hypothesis} as an explicit hypothesis. The
trace identities are exact given the imported Paper~1 K-multiset and
the definitions above.
\end{abstract}

\section{Introduction}\label{sec:intro}

The Hubble tension --- the $\sim 5\sigma$ disagreement between
early-Universe ($\Hzero \approx 67.4$~km/s/Mpc from
Planck~2018~\citep{Planck2020}) and late-Universe
($\Hzero = 73.04 \pm 1.04$~km/s/Mpc from
SH0ES~2022~\citep{Riess2022SH0ES}) inferences of the present-day
Hubble constant --- and the related $\sigeight / S_8$ tension between
Planck and weak-lensing surveys (KiDS-1000~\citep{Heymans2021KiDS},
DES-Y3~\citep{Abbott2022DESY3}) have driven a large literature of
proposed
resolutions~\citep{DiValentinoEtAl2021Tension}.

This paper does \emph{not} propose a cosmological model. Instead, we
report two finite-group facts about the binary icosahedral
group $\twoI = \Vsixhundred$, and observe that
the corresponding rational ratios numerically match published tension
ratios. We organise the paper around three explicit layers of claim:

\begin{itemize}[leftmargin=2em]
  \item \textbf{Layer~1 (Theorem-grade):} exact rational trace
  identities on the $12$-dimensional rational $\Ttau$-cycle space.
  \item \textbf{Layer~2 (Observation):} numerical match between those
  ratios and currently-published tension benchmarks.
  \item \textbf{Layer~3 (Hypothesis):} a coupling rule from the
  operators to ΛCDM observables; \emph{stated as a hypothesis, not
  claimed to be derived}.
\end{itemize}

The reader is invited to evaluate Layers~1 and 2 on their own merits
and to treat Layer~3 with appropriate scepticism.

\section{The $12$-dimensional cycle space and K-class projectors}\label{sec:setup}

We import the following from Paper~1~\citep{V600Paper1}; the present
paper does not re-prove any of these statements.

\begin{itemize}[leftmargin=2em]
\item $\Vsixhundred = \twoI$, the binary icosahedral group of order
$120$, viewed as the vertex set of the regular $600$-cell.
\item A choice of $\tau \in \twoI$ with $\mathrm{ord}(\tau) = 10$
gives the cycle decomposition $T_\tau : g \mapsto \tau g$ of
$\Vsixhundred$ into 12 cycles each of length 10.
\item Each cycle is assigned a non-negative rational K-class via the
K-multiset $\{72{:}1,\,0{:}1,\,52{:}5,\,20{:}5\}$. Concretely, of
the $12$ cycles: one has K-class~$72$, one has K-class~$0$, five
have K-class~$52$, and five have K-class~$20$.
\item $H = \Dicfive \subset \twoI$ is the bulk subgroup formed by the
$K=72$ and $K=0$ cycles ($|H| = 20$).
\end{itemize}

\begin{definition}[Cycle space]\label{def:cyclespace}
Let $\CQ$ be the $12$-dimensional rational vector space with basis
$\{e_i\}_{i=1}^{12}$ indexed by the $12$ $T_\tau$-cycles. Equivalently
$\CQ \cong \mathbb{Q}^{12}$. The basis is determined by the
deterministic cycle ordering supplied by
\texttt{vfd\_v600.group.build\_state()}.
\end{definition}

\begin{definition}[K-class projectors]\label{def:Kproj}
For each $K \in \{72, 0, 52, 20\}$, the \emph{K-class projector} is
the diagonal operator $\PK{K} \in \mathrm{End}_{\mathbb{Q}}(\CQ)$
defined by
\[
  \PK{K}(e_i) \;=\; \begin{cases} e_i & \text{if cycle $i$ has K-class }K, \\
                                  0 & \text{otherwise}. \end{cases}
\]
The identity operator on $\CQ$ is denoted $\IC$.
\end{definition}

The ranks of the four K-class projectors are exactly the K-multiset
multiplicities:
\[
  \mathrm{rk}\,\PK{72} = 1, \quad
  \mathrm{rk}\,\PK{0}  = 1, \quad
  \mathrm{rk}\,\PK{52} = 5, \quad
  \mathrm{rk}\,\PK{20} = 5,
\]
and they sum to $\IC$:
$\PK{72} + \PK{0} + \PK{52} + \PK{20} = \IC$.

\begin{remark}[Baseline rule]\label{rem:baseline}
We define the \emph{ΛCDM-baseline operator} as $\IC$ itself: the
\textbf{K-blind uniform cycle weighting} that assigns coefficient $1$
to every cycle basis vector. Once the baseline is so defined, $\IC$
is unique by construction. We do not claim $\IC$ is the unique
$\tau_\sigma$-invariant rank-$12$ operator on $\CQ$ --- many other
operators commute with the $\tau_\sigma$-induced cycle action; the
baseline is a definitional choice, not a uniqueness theorem.
\end{remark}

\section{The two operators $\Hop$ and $\Cop$}\label{sec:operators}

\begin{definition}[Operators]\label{def:ops}
\[
  \Hop \;:=\; \IC + \PK{72}, \qquad \Cop \;:=\; \IC - \PK{0}.
\]
$\Hop$ is the \emph{max-K-anchor enhancement} of the baseline
($\Hop(e_i) = 2 e_i$ on the unique K-class-$72$ cycle and $e_i$
otherwise). $\Cop$ is the \emph{trivial-K-class suppression} of the
baseline ($\Cop(e_i) = 0$ on the unique K-class-$0$ cycle and $e_i$
otherwise).
\end{definition}

The two operators are diagonal in the cycle basis with entries in
$\{0, 1, 2\}$. Both are self-adjoint with respect to the standard
$\mathbb{Q}$-bilinear form on $\CQ$.

\section{Trace theorem (Layer~1)}\label{sec:trace}

\begin{theorem}[Exact trace identities]\label{thm:trace}
\[
  \mathrm{tr}(\Hop) \;=\; 13, \qquad
  \mathrm{tr}(\Cop) \;=\; 11.
\]
Equivalently,
$\mathrm{tr}(\Hop)/12 = 13/12$ and $\mathrm{tr}(\Cop)/12 = 11/12$,
both exact rationals.
\end{theorem}

\begin{proof}
For any diagonal projector $\PK{K}$, the trace equals its rank, which
equals the K-class multiplicity $m_K$. The K-multiset gives
$m_{72} = 1$ and $m_{0} = 1$; trace is additive on diagonal operators,
so
\[
  \mathrm{tr}(\Hop) = \mathrm{tr}(\IC) + \mathrm{tr}(\PK{72}) = 12 + 1 = 13,
\]
\[
  \mathrm{tr}(\Cop) = \mathrm{tr}(\IC) - \mathrm{tr}(\PK{0}) = 12 - 1 = 11.\qedhere
\]
\end{proof}

\begin{theorem}[K-saturated admissibility]\label{thm:admissibility}
Let $\mathcal{A}_K \subset \mathrm{End}_{\mathbb{Q}}(\CQ)$ denote the
K-saturated diagonal projector algebra, namely the
$\mathbb{Q}$-subalgebra generated by
$\{\IC, \PK{72}, \PK{0}, \PK{52}, \PK{20}\}$.
Within $\mathcal{A}_K$, the only rank-one diagonal projector
corrections to $\IC$ --- that is, the only signed corrections of the
form $\IC \pm Q$ with $Q \in \mathcal{A}_K$ a rank-one projector ---
are $\IC \pm \PK{72}$ and $\IC \pm \PK{0}$.
\end{theorem}

\begin{proof}
The four projectors $\{\PK{72}, \PK{0}, \PK{52}, \PK{20}\}$ are
mutually orthogonal idempotents that resolve the identity, so any
$Q \in \mathcal{A}_K$ has the unique decomposition
\[
  Q \;=\; a_{72}\PK{72} + a_{0}\PK{0} + a_{52}\PK{52} + a_{20}\PK{20},
  \qquad a_K \in \mathbb{Q}.
\]
Idempotence $Q^2 = Q$ on the orthogonal decomposition forces
$a_K^2 = a_K$ for each $K$, hence $a_K \in \{0, 1\}$. The rank of
such a $Q$ is the K-multiplicity sum
$\sum_K a_K\,m_K$ over those $K$ with $a_K = 1$. With the K-multiset
$\{72{:}1,\,0{:}1,\,52{:}5,\,20{:}5\}$, the only $0/1$
coefficient choices giving rank exactly~$1$ are
$(a_{72}, a_0, a_{52}, a_{20}) = (1,0,0,0)$ and
$(0,1,0,0)$, yielding $Q = \PK{72}$ and $Q = \PK{0}$. The signed
rank-one corrections to $\IC$ are therefore exactly
$\pm \PK{72}$ and $\pm \PK{0}$.
\end{proof}

\begin{remark}[Hypotheses are essential]\label{rem:hyp}
Theorem~\ref{thm:admissibility} requires both stated hypotheses:
``diagonal in the cycle basis'' and ``K-saturated''. Dropping either
one introduces a large family of rank-one projectors outside
$\mathcal{A}_K$ (for instance any rank-one orthogonal projector on
$\CQ$ that is not K-saturated, or the rank-one $\tau_\sigma$-invariant
projector onto $\mathrm{span}_{\mathbb{Q}}(e_{52} + e_{20})$ for any
$\tau_\sigma$-paired pair of $K=52, K=20$ basis vectors --- this
projector commutes with the $\tau_\sigma$ cycle action but is not
diagonal in the cycle basis and is not in $\mathcal{A}_K$). The
narrow theorem above is the strongest form we claim.
\end{remark}

\subsection{Phase-independence across the $\mathbb{Z}_2^5$ lifts}

Paper~3~\citep{V600Paper3} produces $32$ canonical $\tau_\sigma$
involutions on $\Vsixhundred$ differing by independent
$\mathbb{Z}_2$ phase choices on each of the five non-bulk right
cosets, and all 32 maps swap the K~$=52$ and K~$=20$ cycle classes
within each coset. We show the present trace identities are
phase-independent.

\begin{lemma}[Phase-independence]\label{lem:phaseind}
For any $\tau_\sigma$ in the $\mathbb{Z}_2^5$ orbit of canonical lifts
of Paper~3, the induced action on the cycle space $\CQ$ exchanges
each pair of (K~$=52$, K~$=20$) cycles within a non-bulk right coset
of $\Dicfive$. In particular the K-class projectors $\PK{72}$ and
$\PK{0}$ commute with the $\tau_\sigma$-induced cycle action for
every choice of phases, and Theorem~\ref{thm:trace} and
Theorem~\ref{thm:admissibility} are independent of phase choice.
\end{lemma}

\begin{proof}
Each per-coset phase swap exchanges a $K=52$ cycle with a $K=20$
cycle within the same right coset; it leaves the $K=72$ and $K=0$
bulk cycles fixed setwise (in fact pointwise). The cycle-level
permutation induced by any $\tau_\sigma$ in the orbit therefore acts
trivially on the basis vectors corresponding to $K \in \{72, 0\}$
and acts within the $K = 52 \leftrightarrow K = 20$ pair on the
other ten basis vectors. In all cases the projectors $\PK{72}$ and
$\PK{0}$ commute with this action, so the traces of operators built
from them are unchanged.
\end{proof}

\section{Numerical observation (Layer~2)}\label{sec:observation}

The exact rationals $13/12$ and $11/12$ from
Theorem~\ref{thm:trace} can be compared numerically to selected
published benchmarks. We tabulate published values, propagate the
trace ratios, and report residuals. \emph{No statistical inference
is claimed; this section is observational only.}

We use the lensing parameter $S_8 := \sigeight \sqrt{\Omega_m / 0.3}$
for the clustering row, since published weak-lensing tension
intervals are reported as $S_8$ rather than bare $\sigeight$.

\begin{table}[h]
\centering
\caption{Layer-2 numerical match. Planck~2018 baseline values from
\citealp{Planck2020}: $\Hzero^\mathrm{Planck} = 67.36$~km/s/Mpc and
$S_8^\mathrm{Planck} = 0.832 \pm 0.013$.}
\label{tab:match}
\begin{tabular}{lllll}
\toprule
Quantity & Trace ratio & Predicted (Layer~2) & Late-Universe value & Residual \\
\midrule
$\Hzero$ [km/s/Mpc] & $13/12 \approx 1.0833$ & $72.97$ & $73.04 \pm 1.04$~\citep{Riess2022SH0ES} & $\approx 0.06\sigma$ \\
$S_8$               & $11/12 \approx 0.9167$ & $0.763$ & $0.766^{+0.020}_{-0.014}$~\citep{Heymans2021KiDS} & $\approx 0.24\sigma$ (KiDS only) \\
\bottomrule
\end{tabular}
\end{table}

The $\Hzero$ row uses the headline SH0ES~2022 value; multiplying the
Planck value by $13/12$ gives $72.97$, which is $\approx 0.06\sigma$
below the SH0ES central value at SH0ES uncertainty $\pm 1.04$
km/s/Mpc.

The $S_8$ row uses the KiDS-1000 multi-probe (3$\times$2pt) headline
$S_8 = 0.766^{+0.020}_{-0.014}$~\citep{Heymans2021KiDS}; multiplying
Planck's $S_8 = 0.832$ by $11/12$ gives $0.763$, which is below the
KiDS-1000 central value. Because the predicted value lies on the
\emph{lower} side of the KiDS asymmetric interval, the directional
residual uses the lower KiDS error $-0.014$:
$(0.766 - 0.763)/0.014 \approx 0.24\sigma$. (Including the Planck
$S_8$ uncertainty $\pm 0.013$ in quadrature with the lower KiDS error
gives $\approx 0.18\sigma$.) The cosmic-shear-only KiDS-1000 result
$S_8 = 0.759^{+0.024}_{-0.021}$~\citep{Asgari2021KiDSCS} places the
prediction $0.763$ inside the central band on the upper side. We do
NOT claim a sub-$0.1\sigma$ match for $S_8$.

\subsection{Ablation: rank-$5$ K-class corrections}

For transparency we record the rank-$5$ alternatives: the trace
ratios that would arise from $\IC + \PK{52}$ or $\IC - \PK{20}$
are $17/12 \approx 1.417$ and $7/12 \approx 0.583$. These rank-$5$
corrections give trace ratios far outside the cited $\Hzero$ and
$S_8$ bands. Within the rank-$1$ pair, the trace data alone
do \emph{not} distinguish $\PK{72}$ from $\PK{0}$ (both have
trace~$1$); the assignment of the two rank-one K-classes to specific
observables is the content of Hypothesis~\ref{hyp:coupling}, not of
the trace identity.

\section{Scope and hypothesis (Layer~3)}\label{sec:hypothesis}

The numerical match in Section~\ref{sec:observation} would be
upgraded to a derivation only by an explicit \emph{coupling rule}
between the cycle-space operators and the cosmological observables.
We state this rule as an explicit hypothesis, and we do NOT claim it
is forced by the trace identity.

\begin{hypothesis}[Cycle-mode coupling, $S_8$ form]\label{hyp:coupling}
There exists a coupling map
$\Gamma : \{\text{scale-rate}, \text{clustering-amplitude}\} \to
\mathcal{A}_K$
such that
\[
  \Gamma(\text{scale-rate}) \;=\; +\PK{72}, \qquad
  \Gamma(\text{clustering-amplitude}) \;=\; -\PK{0},
\]
and the observed late-Universe values of $\Hzero$ and the lensing
parameter $S_8$ relate to the early-Universe ΛCDM values by
\[
  \Hzero^{\mathrm{late}}
  \;=\;
  \frac{\mathrm{tr}(\IC + \Gamma(\text{scale-rate}))}{\mathrm{tr}(\IC)}
  \;\Hzero^{\mathrm{early}},
\qquad
  S_8^{\mathrm{late}}
  \;=\;
  \frac{\mathrm{tr}(\IC + \Gamma(\text{clustering-amplitude}))}{\mathrm{tr}(\IC)}
  \;S_8^{\mathrm{early}}.
\]
We state the hypothesis on $S_8$ rather than on bare $\sigeight$
because published weak-lensing tension intervals are reported as
$S_8$. Translating the second relation to bare $\sigeight$ would
require an additional assumption about how $\Omega_m$ enters the
coupling, which we do not make here.
\end{hypothesis}

We emphasise:

\begin{itemize}[leftmargin=2em]
\item Hypothesis~\ref{hyp:coupling} is \emph{not} a theorem of this
paper. It is a phenomenological postulate that, if assumed, would
make the numerical match of Section~\ref{sec:observation} a
parameter-free consequence of Hypothesis~\ref{hyp:coupling}.
\item The hypothesis is falsifiable: future $\Hzero$ measurements
landing outside the $13/12$-translated band, or $S_8$ measurements
outside the $11/12$-translated band, would refute
Hypothesis~\ref{hyp:coupling} but leave the trace identities
untouched.
\item Mode-counting alone --- e.g.\ counting orthonormal
$\sigma$-antisymmetric modes on $\CQ$ --- does not by itself select
the $13/12$ and $11/12$ ratios. Older notes in this repository
(\texttt{papers/cosmological-folding-rate/dynamics/OD1\_STATUS.md}
and \texttt{od1\_cycle\_spectrum.py}) attempted such a mode-count
and reached different intermediate numbers (e.g.\ $17/12$ from a
different counting). The current paper
abandons mode-counting as a route to the coupling rule and treats
the rule as an explicit hypothesis. A genuine derivation would
likely require an $H_4$-representation decomposition of the
$\sigma$-antisymmetric pair module, which is out of scope for this
paper.
\end{itemize}

\section{Conclusion}\label{sec:conclusion}

\paragraph{What is proved (Layer~1).}
The two operators
$\Hop = \IC + \PK{72}$ and $\Cop = \IC - \PK{0}$
on the $12$-dimensional rational $\Ttau$-cycle space of
$\Vsixhundred$ have exact rational traces $\mathrm{tr}(\Hop) = 13$
and $\mathrm{tr}(\Cop) = 11$, hence ratios $13/12$ and $11/12$ over
the baseline $\mathrm{tr}(\IC) = 12$. Within the K-saturated
diagonal projector algebra, $\pm\PK{72}$ and $\pm\PK{0}$ are the
only rank-one projector corrections to $\IC$. Both results are
phase-independent across the $\mathbb{Z}_2^5$ canonical
$\tau_\sigma$ orbit.

\paragraph{What is observed (Layer~2).}
Multiplying Planck~2018's $\Hzero = 67.36$ by $13/12$ gives
$72.97$~km/s/Mpc, $\approx 0.06\sigma$ below the SH0ES~2022 central
value $73.04 \pm 1.04$. Multiplying Planck~2018's
$S_8 = 0.832 \pm 0.013$ by $11/12$ gives $0.763$. Since the
prediction is below the KiDS-1000 multi-probe central value
$0.766^{+0.020}_{-0.014}$, the directional residual uses the lower
KiDS error $-0.014$:
$(0.766 - 0.763)/0.014 \approx 0.24\sigma$. We do not claim either
match is significant beyond the level reported here.

\paragraph{What is hypothesised (Layer~3).}
Hypothesis~\ref{hyp:coupling} would upgrade Layer~2 to a derivation;
this paper makes no claim to have proved it. Layers 1 and 2 stand
without Layer 3.

\appendix

\section{Verification certificate}\label{app:verify}

The companion script \texttt{verify.py} performs the following exact
checks using only the Python standard library and the shared
\texttt{vfd\_v600} package:
\begin{enumerate}[leftmargin=2em]
\item State construction: $|\Vsixhundred| = 120$, exactly 12
$T_\tau$-cycles each of length $10$, K-multiset
$\{72{:}1,\,0{:}1,\,52{:}5,\,20{:}5\}$.
\item Projector ranks: $\mathrm{rk}\,\PK{72} = \mathrm{rk}\,\PK{0} = 1$
and $\mathrm{rk}\,\PK{52} = \mathrm{rk}\,\PK{20} = 5$.
\item Resolution of identity:
$\PK{72} + \PK{0} + \PK{52} + \PK{20} = \IC$.
\item Trace identities: $\mathrm{tr}(\Hop) = 13$ and
$\mathrm{tr}(\Cop) = 11$, both as exact \texttt{Fraction}s.
\item Trace ratios: $\mathrm{tr}(\Hop)/12 = 13/12$ and
$\mathrm{tr}(\Cop)/12 = 11/12$.
\item Phase-independence (Lemma~\ref{lem:phaseind}): for the
canonical phase-$(0,0,0,0,0)$ $\tau_\sigma$ from Paper~3, AND for
each of the $2^5 = 32$ phase combinations enumerated explicitly,
the induced cycle action exchanges the (K=52, K=20) pairs within
each non-bulk right coset and fixes the K=72 and K=0 bulk cycles;
moreover the diagonal vectors of $\Hop$ and $\Cop$ are pointwise
fixed under the cycle permutation (a stronger statement than
trace invariance).
\item Ablation: the rank-$5$ alternatives $\IC + \PK{52}$ and
$\IC - \PK{20}$ give traces $17$ and $7$, ratios
$17/12 \approx 1.417$ and $7/12 \approx 0.583$.
\end{enumerate}
Reproducibility command:
\begin{verbatim}
PYTHONPATH=papers/v600-programme/lib \
    python3 papers/v600-cosmic-tensions/verify.py
\end{verbatim}

\bibliographystyle{plainnat}
\bibliography{references}

\end{document}
